% SEGMENT OLP-0039-S01
% Part:sets-functions-relations
% Chapter: size-of-sets
% Section: non-enumerability-alt

\documentclass[../../../include/open-logic-section]{subfiles}

\begin{document}

\olfileid{sfr}{siz}{nen-alt}

\olsection{\printtoken{S}{nonenumerable} கணங்கள்}

\begin{editorial}
  இந்தப் பிரிவு \olref[enm-alt]{sec} இன் வரையறைகளைப் பயன்படுத்தி
  $\Bin^\omega$, $\Pow{\Nat}$ ஆகியவற்றின் எண்ணிடத்தகாத தன்மையை
  நிறுவுகிறது; அதாவது $\PosInt$ இலிருந்து வரும் மேற்கோர்த்தலுக்குப்
  பதிலாக $\Nat$ உடனான இருபுறச் சார்பு தேவைப்படுகிறது.
\end{editorial}

% SEGMENT OLP-0039-S02
\begin{explain}
இயல் எண்களின் கணம் $\Nat$ முடிவுறாதது.
அது தெளிவாகவே !!{enumerable} ஆகவும் உள்ளது.
ஆனால் \emph{!!{nonenumerable}} கணங்கள், அதாவது !!{enumerable}
ஆக இல்லாத கணங்கள், இருக்கின்றன என்பது வியப்பூட்டும் உண்மை
(\olref[sfr][siz][enm-alt]{defn:enumerable} ஐப் பார்க்கவும்).

இது வியப்பாகத் தோன்றலாம்.
ஏனெனில் $A$ என்பது !!{nonenumerable} என்று கூறுவது,
!!{bijection} $f \colon \Nat \to A$ \emph{எதுவும் இல்லை} என்று
கூறுவதாகும்; அதாவது, $\Nat$ இன் முடிவுறாத எண்ணிக்கையிலான
!!{element}s[acc] $A$ க்கு அனுப்பும் எந்தச் சார்பும்
$A$ முழுவதையும் தீர்த்துவிடுவதில்லை.
எனவே $A$ என்பது !!{nonenumerable} எனில்,
இயல் எண்களின் எண்ணிக்கையைவிட ``அதிகமான'' !!{element}s
$A$ இல் உள்ளன.

ஒரு கணம் !!{nonenumerable} என நிறுவுவதற்கு,
தகுந்த !!{bijection} எதுவும் இருக்க முடியாது எனக் காட்ட வேண்டும்.
இதற்கான சிறந்த வழி, $A$ இன் !!{element}s[acc] பட்டியலிடும் ஒவ்வொரு
முயற்சியும் குறைந்தது ஓர் !!{element}[acc]யாவது விட்டுவிட வேண்டும்
எனக் காட்டுவதாகும்; $f\colon \Nat \to A$ என்ற எந்தச் சார்பும்
!!{surjective} பண்புடையதல்ல என்பதை இது காட்டுகிறது.
இதனை நிறுவுவதற்கான ஒரு பொதுவான உத்தி,
காண்டோரின் \emph{மூலைவிட்ட முறையைப்} பயன்படுத்துவதாகும்.
$A$ இன் !!{element}s கொண்ட $x_1$, $x_2$, \dots என்ற ஒரு பட்டியல்
தரப்பட்டால், அதன் கட்டுமானத்தினாலேயே அந்தப் பட்டியலில் இடம்பெற
முடியாத $A$ இன் மற்றோர் !!{element}[acc]க் கட்டமைக்கிறோம்.

இவை அனைத்தையும் ஓர் எடுத்துக்காட்டின் மூலம் மிகச் சிறப்பாகப்
புரிந்துகொள்ளலாம்.
எனவே நமது முதல் எடுத்துக்காட்டு, $0$ களாலும் $1$ களாலும் ஆன எல்லா
முடிவுறாத தொடர்களின் கணம் $\Bin^\omega$ ஆகும்.
(`$\Bin$' என்பது இருமத்தைக் குறிக்கிறது; அதை $\{0,1\}$ என்ற
இரண்டு உறுப்புக் கணமாகவே கருதலாம்.)
\oliflabeldef{sfr:card-arithmetic:card-opps:sec}{\footnote{மேலும்
துல்லியமாகச் சொன்னால், $\Bin^\omega$ என்பது $0$ களாலும் $1$s
ஆலும் ஆன எல்லா $\omega$-தொடர்களின் கணம், அதாவது
$\funfromto{\omega}{\{0,1\}}$ என்ற கணம், என வரையறுக்க வேண்டும்.
ஆனால் இதன் பொருள் \olref[sfr][card-arithmetic][card-opps]{sec}
இல் மட்டுமே தெளிவாகும்.}}{}
இருப்பினும், இப்போதைக்கு இந்தச் சற்றுத் தளர்வான சொல்லாக்கம்
குழப்பம் எதையும் ஏற்படுத்தாது.
\end{explain}

% SEGMENT OLP-0039-S03
\begin{thm}
\ollabel{thm:nonenum-bin-omega}
$\Bin^\omega$ என்பது !!{nonenumerable} ஆகும்.
\end{thm}

% SEGMENT OLP-0039-S04
\begin{proof}
$\Bin^\omega$ இன் ஓர் உட்கணத்திற்கான எந்தப் பட்டியலாக்கத்தையும்
கருதுக. எனவே $s_{0}$, $s_{1}$, $s_{2}$, \dots{} என்ற ஏதோ ஒரு
பட்டியல் நமக்கு உள்ளது; ஒவ்வொரு $s_n$ உம் $0$ களாலும் $1$ களாலும்
ஆன ஒரு முடிவுறாத தொடராகும்.
% SOURCE CORRECTION OLSIZ-008; see translation/size-notes.tex.
இந்தப் பட்டியலில் உள்ள $n$ ஆவது தொடரின் $m$ ஆவது இலக்கத்தை
$s_n(m)$ என்க.
இப்போது நமது பட்டியலை ஓர் அட்டவணையாகக் கருதலாம்;
அதில் $s_n(m)$ என்பது $n$ ஆவது கிடைவரிசையிலும்
$m$ ஆவது நெடுவரிசையிலும் வைக்கப்பட்டுள்ளது:
\[
\begin{array}{c|c|c|c|c|c}
& 0 & 1 & 2 & 3 & \dots \\ \hline
0 & \mathbf{s_{0}(0)} & s_{0}(1) & s_{0}(2) & s_0(3) & \dots \\ \hline
1 & s_{1}(0)& \mathbf{s_{1}(1)} & s_1(2) & s_1(3) & \dots \\ \hline
2 & s_{2}(0)& s_{2}(1) & \mathbf{s_2(2)} & s_2(3) & \dots \\ \hline
3 & s_{3}(0)& s_{3}(1) & s_3(2) & \mathbf{s_3(3)} & \dots \\ \hline
\vdots & \vdots & \vdots & \vdots & \vdots & \mathbf{\ddots}
\end{array}
\]
இந்தப் பட்டியலில் இல்லாத, $0$ களாலும் $1$ களாலும் ஆன
$d$ என்ற ஒரு முடிவுறாத தொடரை இப்போது கட்டமைப்போம்.
அதன் ஒவ்வோர் உள்ளீட்டையும் குறிப்பிடுவதன் மூலம்,
அதாவது அனைத்து $n \in \Nat$ க்கும் $d(n)$ ஐக் குறிப்பிடுவதன் மூலம்,
இதனைச் செய்வோம்.
உள்ளுணர்வின்படி, மேலுள்ள அட்டவணையின் மூலைவிட்டத்தில் கீழ்நோக்கிப்
படித்து (இதனால்தான் ``மூலைவிட்ட முறை'' என்ற பெயர்),
% SOURCE CORRECTION OLSIZ-009; see translation/size-notes.tex.
ஒவ்வொரு $1$ ஐயும் $0$ ஆகவும், ஒவ்வொரு $0$ ஐயும் $1$ ஆகவும்
மாற்றுகிறோம்.
மேலும் அருவமாகச் சொன்னால், மூலைவிட்டத்தின் $n$ ஆவது !!{element}
$s_n(n)$ ஆனது $1$ அல்லது $0$ ஆக இருப்பதைப் பொறுத்து,
$d(n)$ ஐ $0$ அல்லது $1$ ஆகப் பின்வருமாறு வரையறுக்கிறோம்:
\[
d(n) =
\begin{cases}
1 & \text{$s_{n}(n) = 0$ எனில்}\\
0 & \text{$s_{n}(n) = 1$ எனில்}
\end{cases}
\]
இந்தத் தொடர் $0$ களாலும் $1$ களாலும் ஆன ஒரு முடிவுறாத தொடர் என்பதால்,
$d \in \Bin^\omega$ என்பது தெளிவு.
ஆனால் எந்த $n \in \Nat$ க்கும் $d(n) \neq s_n(n)$ ஆகுமாறு
$d$ ஐக் கட்டமைத்துள்ளோம்.
அதாவது, $s_n$ இன் $n$ ஆவது உள்ளீட்டிலிருந்து $d$ வேறுபடுகிறது.
எனவே எந்த $n\in \Nat$ க்கும் $d \neq s_n$.
ஆகவே $s_0$, $s_1$, $s_2$, \dots என்ற பட்டியலில் $d$ இருக்க முடியாது.

$\Bin^\omega$ இன் ஏதோ ஓர் உட்கணத்திற்கான எந்தப் பட்டியலாக்கமும்
$\Bin^\omega$ இன் ஏதாவது ஓர் !!{element}[acc] விட்டுவிடும் என்று
காட்டியுள்ளோம்.
எனவே $\Bin^\omega$ என்ற கணத்திற்குப் பட்டியலாக்கம் இல்லை;
அதாவது $\Bin^\omega$ என்பது !!{nonenumerable} ஆகும்.
\end{proof}

% SEGMENT OLP-0039-S05
\begin{explain}
இந்த நிறுவல் முறை, $d$ ஐ வரையறுக்க அட்டவணையின் மூலைவிட்டத்தைப்
பயன்படுத்துவதால் ``மூலைவிட்டமாக்கம்'' எனப்படும்.
ஆயினும், மூலைவிட்டமாக்கத்திற்கு ஓர் அட்டவணை இருக்க வேண்டியதில்லை.
உண்மையில், அட்டவணையோ உண்மையான மூலைவிட்டமோ இல்லாதபோதும்,
இதே போன்ற கருத்தைப் பயன்படுத்தி ஒரு கணம் !!{nonenumerable}
எனக் காட்டலாம். பின்வரும் முடிவு இதனை விளக்குகிறது.
\end{explain}

% SEGMENT OLP-0039-S06
\begin{thm}
\ollabel{thm:nonenum-pownat}
$\Pow{\Nat}$ என்பது !!{enumerable} அல்ல.
\end{thm}

% SEGMENT OLP-0039-S07
\begin{proof}
$\Nat$ இன் உட்கணங்களைக் கொண்ட ஒவ்வொரு பட்டியலும்
$\Nat$ இன் ஏதாவது ஓர் உட்கணத்தை விட்டுவிடுகிறது எனக் காட்டி,
அதே முறையில் தொடர்கிறோம்.
எனவே $\Nat$ இன் உட்கணங்களைக் கொண்ட
$N_0, N_1, N_2, \ldots$ என்ற ஏதோ ஒரு பட்டியல் நமக்கு உள்ளது எனக் கொள்க.
$n \in D$ ஆக இருக்கும்போது, அப்போதும் அப்போது மட்டுமே,
$n \notin N_{n}$ ஆகும் வகையில் $D$ என்ற கணத்தைப் பின்வருமாறு
வரையறுக்கிறோம்:
\[
D = \Setabs{n \in \Nat}{n \notin N_n}
\]
$D\subseteq \Nat$ என்பது தெளிவு.
ஆனால் $D$ பட்டியலில் இருக்க முடியாது.
ஏனெனில் கட்டுமானத்தின்படி, $n \in D$ ஆக இருக்கும்போது,
அப்போதும் அப்போது மட்டுமே, $n\notin N_n$;
எனவே எந்த $n \in \Nat$ க்கும் $D \neq N_n$.
\end{proof}

% SEGMENT OLP-0039-S08
\begin{explain}
முந்தைய நிறுவலில் மூலைவிட்டம் குறிப்பிடப்படவில்லை.
இருப்பினும், அதில் ஒரு மூலைவிட்டம் இருப்பதாகப் பின்வருமாறு
காட்சிப்படுத்திக் கொள்ளலாம்:
$N_0$, $N_1$, \dots என்ற கணங்கள் ஓர் அட்டவணையில் எழுதப்பட்டுள்ளதாகக்
கற்பனை செய்க; $m \in N_n$ ஆக இருக்கும்போது, அப்போதும் அப்போது மட்டுமே,
$n$ ஆவது கிடைவரிசையின் $m$ ஆவது நெடுவரிசையில் $m$ ஐ எழுதி,
$N_n$ ஐ அந்தக் கிடைவரிசையில் எழுதுகிறோம்.
எடுத்துக்காட்டாக, அந்தப் பட்டியலின் முதல் நான்கு கணங்கள்
$\{0,1,2,\dots\}$, $\{1, 3, 5, \dots\}$, $\{0,1,4\}$,
$\{2,3,4,\dots\}$ என்க; அப்போது நமது அட்டவணை பின்வருமாறு தொடங்கும்:
\[
\begin{array}{r@{}rrrrrrr}
  N_0 = \{ & \mathbf{0}, & 1, & 2, & & & & \dots\}\\
  N_1 = \{ &  & \mathbf{1}, &  & 3, &  & 5, & \dots\}\\
  N_2 = \{ & 0, & 1, &  &  & 4\phantom{,} &  & \}\\
  N_3 = \{ &  &  & 2, & \mathbf{3}, & 4, & & \dots\}\\
  &\vdots & & & & & & \ddots\phantom{\}}
  \end{array}
\]
பின்னர், மூலைவிட்டத்தில் கீழ்நோக்கிச் சென்று,
$n$ மூலைவிட்டத்தில் \emph{இல்லாதபோது}, அப்போதும் அப்போது மட்டுமே,
$n \in D$ ஆகுமாறு பெறப்படும் கணமே $D$.
எனவே மேலுள்ள நிலையில், $0$, $1$ ஆகியவற்றை விட்டுவிடுவோம்;
$2$ ஐச் சேர்ப்போம்; $3$ ஐ விட்டுவிடுவோம்; இவ்வாறே தொடர்வோம்.
\end{explain}

% SEGMENT OLP-0039-S09
\begin{prob}
$f \colon \Nat \to \Nat$ என்ற எல்லாச் சார்புகளின் கணமும்
!!{nonenumerable} என்பதை வெளிப்படையான மூலைவிட்ட வாதத்தால் காட்டுக.
அதாவது, $f_1$, $f_2$, \dots என்பது சார்புகளின் ஒரு பட்டியலாகவும்,
ஒவ்வொரு $f_i\colon \Nat \to \Nat$ எனவும் இருந்தால்,
அந்தப் பட்டியலில் இல்லாத ஏதோ ஒரு $g \colon \Nat \to \Nat$
உள்ளது எனக் காட்டுக.
\end{prob}

\end{document}
